\documentclass[a4paper,10.5pt]{article}

\usepackage[T1]{fontenc}
\usepackage[utf8]{inputenc}
\usepackage[brazil]{babel}
\usepackage[scaled=0.92]{helvet}
\renewcommand{\familydefault}{\sfdefault}
\usepackage[left=1cm,right=1cm,top=0.8cm,bottom=0.5cm]{geometry}
\usepackage{enumitem}
\usepackage{titlesec}
\usepackage{xcolor}
\usepackage{tabularx}
\usepackage{graphicx}
\usepackage{tikz}
\graphicspath{{img/}}
\usepackage[hidelinks]{hyperref}

%--------------------------------------------------------------
% Cores e estilos
%--------------------------------------------------------------
\definecolor{primary}{HTML}{1F3A5F}
\definecolor{muted}{HTML}{444444}

\hypersetup{colorlinks=true, urlcolor=primary}

\titleformat{\section}
  {\normalfont\large\bfseries\color{primary}}
  {}{0em}{}[\vspace{-0.7em}\color{primary}\rule{\textwidth}{0.8pt}]
\titlespacing*{\section}{0pt}{1.0em}{0.55em}

\setlist[itemize]{leftmargin=1.1em, topsep=2pt, itemsep=1.5pt, parsep=0pt}
\pagestyle{empty}
\setlength{\parindent}{0pt}

% Cabeçalho de experiência: {cargo}{empresa}{período}
\newcommand{\cargo}[3]{%
  \vspace{2pt}
  \begin{tabularx}{\textwidth}{@{}Xr@{}}
    \textbf{#1} \textnormal{\textbar{} #2} & \textcolor{muted}{#3}
  \end{tabularx}%
  \vspace{1pt}
}

% Linha de competência: {rótulo}{conteúdo}
\newcommand{\skillline}[2]{\textbf{#1:} #2\\}

% Foto de perfil circular: [diâmetro]{arquivo}
\newcommand{\fotoperfil}[2][2.6cm]{%
  \begin{tikzpicture}
    \begin{scope}
      \clip (0,0) circle (#1/2);
      \node[inner sep=0pt] at (0,0) {\includegraphics[width=#1]{#2}};
    \end{scope}
    \draw[primary, line width=0.9pt] (0,0) circle (#1/2);
  \end{tikzpicture}%
}

\begin{document}

%==============================================================
% CABEÇALHO
%==============================================================
\begin{minipage}[c]{2.6cm}
  \fotoperfil[2.5cm]{foto.jpg}
\end{minipage}%
\hfill
\begin{minipage}[c]{\dimexpr\textwidth-3.0cm\relax}
  {\LARGE\bfseries\color{primary} ROBSON THIAGO DUARTE DA SILVA}\\[3pt]
  {\normalsize Estudante de Engenharia da Computação (UFPE) \textbar{} Desenvolvedor \& Automação de Processos}\\[5pt]
  \small
  Recife -- PE \textbar{}
  (81) 98202-1412 \textbar{}
  \href{mailto:robgol.fidel@gmail.com}{robgol.fidel@gmail.com} \textbar{}
  \href{https://sites.google.com/cin.ufpe.br/rtds/p\%C3\%A1gina-inicial}{Portfólio}
\end{minipage}

\vspace{-0.3em}

%==============================================================
% RESUMO
%==============================================================
\section{Resumo Profissional}
Estudante de Engenharia da Computação na UFPE, com atuação profissional em desenvolvimento
de software e forte vivência em \textbf{automação de processos, integração de sistemas e manipulação de dados}.
Experiência prática construindo automações em \textbf{Python (Selenium, Pandas)} e \textbf{Microsoft Power Automate},
consumindo \textbf{APIs REST} e trabalhando com \textbf{SQL} e bancos de dados para garantir consistência e
interoperabilidade entre plataformas. Vivência com suporte de TI sob framework \textbf{ITIL}, o que consolidou uma
leitura crítica de processos operacionais --- base direta para mapear gargalos e propor soluções de RPA.
Perfil de iniciativa comprovado por \textbf{6 premiações em hackathons} e por um produto homenageado no
\textbf{Prêmio Innovare (CNJ/STF)}.

%==============================================================
% COMPETÊNCIAS
%==============================================================


\section{Competências Técnicas}
\skillline{Automação \& RPA}{Python com Selenium (automação web e de tarefas repetitivas), Pandas (ETL e tratamento de dados), Microsoft Power Automate (fluxos e integrações), automação de planilhas via Google Sheets}
\skillline{Dados \& Integração}{SQL, modelagem e consulta a bancos de dados relacionais, consumo e integração de APIs REST, extração, transformação e análise de dados (ETL)}
\skillline{Desenvolvimento}{React, TypeScript, JavaScript, Python, HTML, CSS}
\skillline{Processos \& Qualidade}{Levantamento e análise de processos, ITIL (suporte N1), testes e validação de soluções, documentação técnica, LGPD}
\skillline{Ferramentas}{Git/GitHub, Google Workspace, ambientes ágeis}
\skillline{Workflow IA}{Claude Code CLI, Codex CLI e Cursor IDE}


%==============================================================
% EXPERIÊNCIA
%==============================================================
\vspace{-0.5em}
\section{Experiência Profissional}

\cargo{Desenvolvedor Front-End}{V-Lab UFPE}{1 ano e 5 meses}
\begin{itemize}
  \item Desenvolvimento de produtos digitais para o \textbf{AVAMEC}, plataforma educacional do Ministério da Educação, com React e TypeScript.
  \item Integração de interfaces com \textbf{APIs REST}, garantindo consistência dos dados entre front-end e serviços de back-end.
  \item Participação no ciclo completo de entrega: análise de requisitos, desenvolvimento, testes, validação e melhoria contínua das soluções.
  \item Colaboração em equipe multidisciplinar, com versionamento em Git e rotina de code review.
\end{itemize}

\cargo{Tutor de Tecnologia}{CESAR School --- Projeto Florescendo Talentos}{2 anos e 3 meses}
\begin{itemize}
  \item Capacitação de estudantes de Ensino Médio e EJA em \textbf{lógica de programação (Python)}, desenvolvimento web (HTML, CSS, JS) e \textbf{manipulação de dados com SQL e Google Sheets}.
  \item Construção de materiais e exercícios práticos de \textbf{automação de planilhas e consultas a dados}, traduzindo conceitos técnicos para públicos não técnicos.
  \item Desenvolvimento de comunicação funcional e didática --- habilidade aplicável ao levantamento de requisitos junto a áreas de negócio.
\end{itemize}

\cargo{Analista de Help Desk Júnior}{Datamétrica}{1 ano}
\begin{itemize}
  \item Suporte de TI Nível 1 a órgãos do Governo do Estado de Pernambuco, aplicando o framework \textbf{ITIL} na gestão de incidentes e requisições.
  \item Mapeamento de \textbf{processos operacionais recorrentes} e identificação de gargalos passíveis de padronização e automação.
  \item Registro, categorização e análise de chamados, gerando base para priorização de melhorias.
\end{itemize}

\cargo{Jovem Aprendiz --- Administrativo}{Proxxi Tecnologia}{1 ano e 6 meses}
\begin{itemize}
  \item Organização e controle de documentos com adequação aos requisitos da \textbf{LGPD}.
  \item Suporte a setores internos, com padronização de rotinas administrativas e controle de informações em planilhas.
\end{itemize}

\cargo{Professor de Informática Básica}{News Center Cursos \textbar{} ONG Coque Connecta (voluntário)}{10 meses}
\begin{itemize}
  \item Ensino de informática básica e ferramentas de produtividade para turmas regulares e para crianças da comunidade do Coque (Recife).
\end{itemize}

%==============================================================
% PROJETOS
%==============================================================
\vspace{-0.5em}
\section{Projetos e Reconhecimentos}

\cargo{Desenvolvedor Full Stack --- Fênix Connect}{Homenageado no Prêmio Innovare (21ª ed.)}{}
\begin{itemize}
  \item Aplicativo voltado à reintegração social de egressos do sistema prisional, homenageado no \textbf{Prêmio Innovare} --- iniciativa do Instituto Innovare em parceria com \textbf{CNJ, STF e Grupo Globo}.
  \item Atuação em front-end e back-end, incluindo modelagem de dados, integração via API e testes de validação.
\end{itemize}

\textbf{Premiações em Hackathons}
\begin{itemize}
    \item \textbf{Encubação} -- Recife Resolve 2026.1 \textbar{} \textbf{2º lugar} -- Hackathon BBTS 2025 \textbar{} \textbf{2º lugar} -- EDscript Hackathon do EDS (2026) \textbar{} \textbf{1º lugar} -- Hacker Cidadão 9.0 \textbar{} \textbf{1º lugar} -- Hackathon de SUAPE (2024) \textbar{} \textbf{2º lugar} -- Hacker Cidadão 7.0 \textbar{} \textbf{2º lugar} -- Ideatel 2025 (UFRPE)
\end{itemize}

%==============================================================
% FORMAÇÃO
%==============================================================
\section{Formação Acadêmica}
\cargo{Bacharelado em Engenharia da Computação}{Universidade Federal de Pernambuco (UFPE)}{Em andamento}
\vspace{2pt}

%==============================================================
% VOLUNTARIADO
%==============================================================
\vspace{-1.5em}
\section{Voluntariado e Comunidade Tech}
\begin{itemize}
  \item \textbf{Rec'n'Play} --- Arena IA e Arena de Inovação (2024 e 2025); \textbf{NASA Space Apps} (2025); \textbf{Secomp UFPE} (2024); \textbf{Agile Trends NE} (2025, diretor de palco); \textbf{Casa do Julgamento} (2025).
\end{itemize}

%==============================================================
% IDIOMAS
%==============================================================
\section{Idiomas}
Português (nativo) \textbar{} Inglês (leitura técnica)

\end{document}
